\subsection{INVERSE OF A CORRESPONDENCE}

Let \(G\) be a graph and \(A = \operatorname{pr}_1 G\) , \(B = \operatorname{pr}_2 G\) its projections. The relation \((y, x) \in G\) implies \((x, y) \in B \times A\) ; this relation therefore has a graph which consists of all ordered pairs \((x, y)\) such that \((y, x) \in G\) .

DEFINITION 5. Let G be a graph. The graph whose elements are the ordered pairs \((x, y)\) such that \((y, x) \in G\) is called the inverse of G and is denoted by \(\overline{G}\) .

For every set X, \(\overline{G}\langle X\rangle\) is called the inverse image of X under G.

It is clear that the inverse of \(\overline{G}\) is \(G\) , and that

\[
\mathrm{pr} _ {1} ^ {- 1} \mathrm{G} = \mathrm{pr} _ {2} \mathrm{G}, \quad \mathrm{pr} _ {2} ^ {- 1} \mathrm{G} = \mathrm{pr} _ {1} \mathrm{G}.
\]

In particular, if X, Y are two sets, we have

\[
\widehat {\mathbf {X} \times \mathbf {Y}} = \mathbf {Y} \times \mathbf {X}.
\]

A graph G is said to be symmetric if \(\overline{G}^{-1}=G\) .

Let \(\Gamma = (G, A, B)\) be a correspondence between \(A\) and \(B\) . Since \(\operatorname{pr}_1^{\overline{-1}} G \subset B\) and \(\operatorname{pr}_2^{\overline{-1}} G \subset A\) , the triple \((G, B, A)\) is a correspondence between \(B\) and \(A\) , called the inverse of the correspondence \(\Gamma\) , and denoted by \(\overline{\Gamma}\) . For every subset \(Y\) of \(B\) , the image \(\overline{\Gamma} \langle Y \rangle\) of \(Y\) under \(\overline{\Gamma}\) is also called the inverse image of \(Y\) under \(\Gamma\) . Clearly the inverse of \(\overline{\Gamma}\) is \(\Gamma\) .

